\documentclass[11pt]{article}
\usepackage[a4paper,margin=1in]{geometry}
\usepackage{fontspec}
\usepackage[boldfont=false]{xeCJK}
\setCJKmainfont{FandolSong-Regular.otf}
\usepackage{amsmath}
\usepackage{amssymb}
\usepackage{array}
\usepackage{booktabs}
\usepackage{graphicx}
\usepackage{adjustbox}
\usepackage{enumitem}
\setlist[itemize]{topsep=2pt,partopsep=0pt,parsep=0pt,itemsep=2pt,leftmargin=1.4em}
\setlist[enumerate]{topsep=2pt,partopsep=0pt,parsep=0pt,itemsep=2pt,leftmargin=1.6em}
\usepackage{parskip}
\usepackage{fvextra}
\fvset{breaklines=true,breakanywhere=true,fontsize=\small}
\setlength{\parindent}{0pt}

\begin{document}

\begin{center}
  {\LARGE\bfseries Operations management}\\[0.35em]
  {\large A Level Business}
\end{center}
\vspace{0.6em}

\section*{What operations management does}

\textbf{Operations} 运营 is the part of a business that makes the product. \textbf{Operations management} 运营管理 is the job of turning \textbf{inputs} 投入 into \textbf{outputs} 产出 well.

This is the \textbf{transformation process} 转化过程:

\begin{itemize}
\item \textbf{inputs} — raw materials, labour, machines and money.

\item \textbf{process} — the work that changes the inputs.

\item \textbf{outputs} — finished goods and services for the customer.

\end{itemize}

Good operations \textbf{add value} 增值: the output is worth more than the inputs used to make it. Using \textbf{raw materials} 原材料 (the basic materials a business starts with) well is a big part of this.

\section*{Productivity and efficiency}

\textbf{Productivity} 生产率 measures how much output you get from your inputs:

\[
\text{productivity} = \frac{\text{output}}{\text{input used}}
\]

\textbf{Efficiency} 效率 means making products with as little waste of time, materials and money as possible. Higher productivity and efficiency lower the cost of each unit, so the firm can charge less or earn more.

\section*{Capital- and labour-intensive operations}

\begin{itemize}
\item a \textbf{capital-intensive} 资本密集型 operation uses mostly machines (e.g. a car factory). It has high set-up costs but low cost per unit, and is good for large amounts.

\item a \textbf{labour-intensive} 劳动密集型 operation uses mostly people (e.g. a hair salon). It is flexible and good for personal service, but quality can vary.

\end{itemize}

\section*{Methods of production}

A firm chooses how to organise production to suit its product and demand.

\begin{center}
\adjustbox{max width=\linewidth}{%
\begin{tabular}{lll}
\toprule
\textbf{Method} & \textbf{What it is} & \textbf{Best for} \\
\midrule
\textbf{job production} 单件生产 & making one unique item at a time & a wedding cake, a bridge \\
\textbf{batch production} 批量生产 & making a group of the same item, then switching & bread in flavours \\
\textbf{flow production} 流水线生产 & making items non-stop on a line & bottled drinks, cars \\
\textbf{mass customisation} 大规模定制 & a flow line that still lets customers choose options & cars with chosen colours \\
\bottomrule
\end{tabular}}
\end{center}

Job production is flexible but costly per unit. Flow production has a low unit cost but needs high, steady demand.

\section*{Sustainability and technology}

\textbf{Sustainability} 可持续性 means meeting today's needs without harming the future or the planet. Sustainable operations cut waste, use less energy, recycle, and use materials that can be replaced.

\textbf{Technology} 技术 is widely used in operations. \textbf{Automation} 自动化 (machines and robots doing tasks) raises output, improves quality and lowers cost, but it has a high set-up cost and can replace jobs.

\section*{Inventory and why it matters}

\textbf{Inventory} 库存 (stock) is the goods a business holds: raw materials, part-finished goods, and finished goods waiting to be sold.

A firm holds inventory so it can keep producing and meet customer orders quickly. But holding inventory has costs:

\begin{itemize}
\item \textbf{holding costs} 持有成本 — storage, insurance, and money tied up in stock.

\item the risk that goods go out of date or are damaged.

\end{itemize}

Holding too little stock is also risky: the firm may run out and lose sales.

\section*{Inventory control charts}

An \textbf{inventory control chart} 库存控制图 shows how stock rises when an order arrives and falls as it is used. Key terms:

\begin{itemize}
\item \textbf{re-order level} 再订货水平 — the stock level at which a new order is placed.

\item \textbf{lead time} 交货周期 — the time between placing an order and receiving it.

\item \textbf{buffer inventory} 缓冲库存 — a minimum "safety" stock kept in case of delays or extra demand.

\item \textbf{maximum inventory} 最大库存 — the most stock the firm wants to hold.

\end{itemize}

\section*{Just-in-time and just-in-case}

There are two main ways to manage stock.

\begin{itemize}
\item \textbf{just-in-time} 准时制 (JIT) — stock arrives only as it is needed, so almost no inventory is held. This cuts holding costs and waste, but a late delivery can stop production.

\item \textbf{just-in-case} 备货制 — extra stock is held in case of problems. This is safer but costs more to store.

\end{itemize}

\section*{Capacity utilisation}

\textbf{Capacity} 产能 is the most a business can produce with its current resources. \textbf{Capacity utilisation} 产能利用率 shows how much of that capacity is being used:

\[
\text{capacity utilisation} = \frac{\text{actual output}}{\text{maximum possible output}} \times 100\%
\]

\begin{itemize}
\item \textbf{under-utilisation} 产能利用不足 (a low figure) means resources sit idle, so the fixed cost spread over each unit is high.

\item \textbf{over-utilisation} 产能过度利用 (very close to 100\%) leaves no time for repairs, can tire staff, and may lower quality.

\end{itemize}

Most firms aim for a high but not full level, around 85–90\%.

\section*{Outsourcing}

\textbf{Outsourcing} 外包 means paying another firm (a \textbf{supplier} 供应商) to do work the business used to do itself, such as cleaning, IT or making parts.

Benefits: it can be cheaper, lets the firm focus on what it does best, and adds flexibility. Drawbacks: less control over quality, possible delays, and the risk of sharing secrets with another firm.


\end{document}
